% {{CONTEST_UPPER}} {{PROBLEM_UPPER}} - {{TITLE}}
% ビルド: ../../tools/build_slides.sh      （latexmk -pdfdvi → platex + dvipdfmx）
% 図は TikZ で描く。書き方の指針は tools/slides_guide.md を参照。
\documentclass[dvipdfmx,aspectratio=169]{beamer}
\usepackage{amsmath}
\usepackage{booktabs}
\usepackage{pxjahyper}
\usepackage{tikz}
\usetikzlibrary{arrows.meta,positioning,backgrounds,fit}

\usetheme{Madrid}
\usecolortheme{default}
\setbeamertemplate{navigation symbols}{}

% 状態の色分け（正解/不正解・採用/棄却などに使い回す）
\definecolor{ok}{RGB}{0,140,70}
\definecolor{ng}{RGB}{200,30,30}
\definecolor{onc}{RGB}{240,180,0}

\title[{{CONTEST_UPPER}} {{PROBLEM_UPPER}}]{AtCoder {{CONTEST_UPPER}} -- {{PROBLEM_UPPER}}\\\small {{TITLE}} の解説}
\author[{{CONTEST}}]{}
\institute{}
\date{}

\begin{document}

%====================================================================
\begin{frame}
  \titlepage
\end{frame}

%====================================================================
\begin{frame}{問題の概要}
  % 入力・出力・制約を 5 行以内で。問題文を丸写ししない。
  % 具体例を 1 つ図にすると、以降の説明が全部それを指せる。
\end{frame}

%====================================================================
\begin{frame}{まず制約を読む}
  % 「なぜこの解法に行き着くか」の起点。計算量の見積りを数式で置く。
\end{frame}

%====================================================================
\begin{frame}{アイデア}
  % 核心をひとつだけ。図（TikZ）と文章を左右 columns に分けると収まりやすい。
  \begin{columns}
    \begin{column}{0.46\textwidth}
      \centering
      \begin{tikzpicture}
      \end{tikzpicture}

      % 図の注釈は tikzpicture の外に置く（中に置くと段の幅をはみ出しやすい）
    \end{column}
    \begin{column}{0.5\textwidth}
    \end{column}
  \end{columns}
\end{frame}

%====================================================================
\begin{frame}[fragile]{擬似コードと実装の要点}
  % verbatim / \verb を使うフレームには [fragile] が必須
  \begin{verbatim}
  \end{verbatim}
\end{frame}

%====================================================================
\begin{frame}{サンプルで確認}
  % 小さい入力を手で全部追う。explanation.md の表をそのまま持ってくる。
\end{frame}

%====================================================================
\begin{frame}[fragile]{まとめ}
  % 次に同じ形を見たとき効く教訓を 3〜4 個 + 同系統の問題
\end{frame}

\end{document}
